criteria for @@M0@@ to be injective (resp. surjective).

\sgaStableTarget{sga6.e10.para.7-18-2}\textbf{7.18.2.} Suppose that we are in one of the following two cases:
\begin{enumerate}[label=\alph*)]
\item The field @@M1@@ is separably closed, and @@M2@@ is a separable extension.
\item The field @@M3@@ is infinite, and @@M4@@ is a pure extension of it.
\end{enumerate}
Then the map @@M5@@ is injective.

Indeed, using the fact that @@M6@@ is the filtered inductive limit of subalgebras of finite type over @@M7@@, with flat transition morphisms, and the commutation of @@M8@@ with filtered inductive limits, one is reduced in case a) to proving that, if @@M9@@ is an integral algebra of finite type over @@M10@@ whose fraction field @@M11@@ is separable over @@M12@@, then @@M13@@ is injective. Now the hypotheses on @@M14@@ and on @@M15@@ imply that there exists a @@M16@@-homomorphism @@M17@@; the composite @@M18@@ is then the identity, whence the conclusion. In case b), passage to the limit similarly reduces the question to proving that, if @@M19@@ is an integral algebra of finite type over @@M20@@ whose fraction field is a pure extension of @@M21@@, then @@M22@@ is injective. Since @@M23@@ is infinite, the hypothesis again implies that there exists a @@M24@@-homomorphism @@M25@@, and one concludes as above.

\sgaStableTarget{sga6.e10.para.7-18-3}\textbf{7.18.3.} Now suppose that @@M26@@ is a functor with values in the category @@M27@@ of abelian groups, and suppose that for every finite extension @@M28@@ of @@M29@@ a norm homomorphism
@@M30@@
has been defined such that the composite
@@M31@@
is multiplication by @@M32@@. Under these conditions:
\begin{enumerate}[label=(\roman*)]
